Аннотация

\hl{Техническое задание устанавливает требования к разработке
программного модуля для сетевого симулятора Ns-3, предназначенного для
моделирования программно-аппаратного гидроакустического
приёмопередатчика и гидроакустического канала связи с целью проведения
виртуальных и полунатурных испытаний и исследований сценариев
эксплуатации. Техническое задание включает в себя страниц~--~38,
приложений~--~5.}

\hl{Список ключевых слов: моделирование, виртуальные испытания,
полунатурные испытания, сетевой симулятор, Ns-3, гидроакустический
приёмопередатчик, гидроакустический канал связи.\\
}

\section{Содержание}\label{ux441ux43eux434ux435ux440ux436ux430ux43dux438ux435}

\hyperref[ux432ux432ux435ux434ux435ux43dux438ux435]{Содержание 4}

\hyperref[ux432ux432ux435ux434ux435ux43dux438ux435]{Введение 6}

\hyperref[ux43dux430ux438ux43cux435ux43dux43eux432ux430ux43dux438ux435-ux43fux440ux43eux433ux440ux430ux43cux43cux44b]{1.1
Наименование программы 6}

\hyperref[ux43eux431ux43bux430ux441ux442ux44c-ux43fux440ux438ux43cux435ux43dux435ux43dux438ux44f]{1.2
Область применения 6}

\hyperref[ux43eux441ux43dux43eux432ux430ux43dux438ux435-ux434ux43bux44f-ux440ux430ux437ux440ux430ux431ux43eux442ux43aux438]{2.
Основание для разработки 6}

\hyperref[ux43dux430ux437ux43dux430ux447ux435ux43dux438ux435-ux440ux430ux437ux440ux430ux431ux43eux442ux43aux438]{3.
Назначение разработки 7}

\hyperref[ux442ux440ux435ux431ux43eux432ux430ux43dux438ux44f-ux43a-ux43fux440ux43eux433ux440ux430ux43cux43cux435]{4
Требования к программе 7}

\hyperref[ux442ux440ux435ux431ux43eux432ux430ux43dux438ux44f-ux43a-ux444ux443ux43dux43aux446ux438ux43eux43dux430ux43bux44cux43dux44bux43c-ux445ux430ux440ux430ux43aux442ux435ux440ux438ux441ux442ux438ux43aux430ux43c]{4.1
Требования к функциональным характеристикам 7}

\hyperref[ux442ux440ux435ux431ux43eux432ux430ux43dux438ux44f-ux43a-ux430ux440ux445ux438ux442ux435ux43aux442ux443ux440ux435-ux43cux43eux434ux435ux43bux438-ux43fux440ux43eux433ux440ux430ux43cux43cux43dux43e-ux430ux43fux43fux430ux440ux430ux442ux43dux43eux433ux43e-ux43fux440ux438ux451ux43cux43eux43fux435ux440ux435ux434ux430ux442ux447ux438ux43aux430]{4.1.2
Требования к архитектуре модели программно-аппаратного приёмопередатчика
8}

\hyperref[ux43eux431ux449ux438ux435-ux442ux440ux435ux431ux43eux432ux430ux43dux438ux44f]{4.1.2.1
Общие требования 8}

\hyperref[ux43dux435ux437ux430ux432ux438ux441ux438ux43cux43eux441ux442ux44c-ux443ux440ux43eux432ux43dux435ux439]{4.1.2.2
Независимость уровней 8}

\hyperref[ux444ux443ux43dux43aux446ux438ux438-ux444ux438ux437ux438ux447ux435ux441ux43aux43eux433ux43e-ux443ux440ux43eux432ux43dux44f]{4.1.2.2
Функции физического уровня 9}

\hyperref[ux442ux440ux435ux431ux43eux432ux430ux43dux438ux44f-ux43a-ux430ux440ux445ux438ux442ux435ux43aux442ux443ux440ux435-ux43cux43eux434ux435ux43bux438-ux433ux438ux434ux440ux43eux430ux43aux443ux441ux442ux438ux447ux435ux441ux43aux43eux433ux43e-ux43aux430ux43dux430ux43bux430-ux441ux432ux44fux437ux438]{4.1.3
Требования к архитектуре модели гидроакустического канала связи 9}

\hyperref[ux442ux440ux435ux431ux43eux432ux430ux43dux438ux44f-ux43a-ux440ux435ux430ux43bux438ux437ux430ux446ux438ux438-ux43cux43eux434ux435ux43bux438-ux43fux440ux438ux451ux43cux43eux43fux435ux440ux435ux434ux430ux442ux447ux438ux43aux430]{4.1.4
Требования к реализации модели приёмопередатчика 10}

\hyperref[ux43eux431ux449ux438ux435-ux43fux43eux43bux43eux436ux435ux43dux438ux44f]{4.1.4.1
Общие положения 10}

\hyperref[ux441ux43eux441ux442ux43eux44fux43dux438ux44f-ux444ux438ux437ux438ux447ux435ux441ux43aux43eux433ux43e-ux443ux440ux43eux432ux43dux44f]{4.1.4.2
Состояния физического уровня 10}

\hyperref[_rgwhsh7we1q0]{4.1.6 Требования к модели гидроакустического
канала связи 12}

\hyperref[_ytskkr1fx40n]{4.1.6.1 Расчёт характеристик распространения
12}

\hyperref[_lhvuj3g4c8ca]{4.1.6.3 Рассчёт затухания (потерь) сигнала 13}

\hyperref[_uhesppck4kdg]{4.2 Требования к надежности 13}

\hyperref[ux442ux440ux435ux431ux43eux432ux430ux43dux438ux44f-ux43a-ux43eux431ux435ux441ux43fux435ux447ux435ux43dux438ux44e-ux43dux430ux434ux451ux436ux43dux43eux433ux43e-ux444ux443ux43dux43aux446ux438ux43eux43dux438ux440ux43eux432ux430ux43dux438ux44f-ux43fux440ux43eux433ux440ux430ux43cux43cux44b]{4.2.1
Требования к обеспечению}
\hyperref[ux442ux440ux435ux431ux43eux432ux430ux43dux438ux44f-ux43a-ux43eux431ux435ux441ux43fux435ux447ux435ux43dux438ux44e-ux43dux430ux434ux451ux436ux43dux43eux433ux43e-ux444ux443ux43dux43aux446ux438ux43eux43dux438ux440ux43eux432ux430ux43dux438ux44f-ux43fux440ux43eux433ux440ux430ux43cux43cux44b]{надёжного
функционирования программы 13}

\hyperref[ux432ux440ux435ux43cux44f-ux432ux43eux441ux441ux442ux430ux43dux43eux432ux43bux435ux43dux438ux44f-ux43fux43eux441ux43bux435-ux43eux442ux43aux430ux437ux430]{4.2.2
Время восстановления после отказа 13}

\hyperref[ux43eux442ux43aux430ux437ux44b-ux438ux437-ux437ux430-ux43dux435ux43aux43eux440ux440ux435ux43aux442ux43dux44bux445-ux434ux435ux439ux441ux442ux432ux438ux439-ux43fux43eux43bux44cux437ux43eux432ux430ux442ux435ux43bux44f]{4.2.3
Отказы из-за некорректных действий пользователя 13}

\hyperref[ux442ux440ux435ux431ux43eux432ux430ux43dux438ux44f-ux43a-ux443ux441ux43bux43eux432ux438ux44fux43c-ux44dux43aux441ux43fux43bux443ux430ux442ux430ux446ux438ux438]{4.3
Требования к условиям эксплуатации 14}

\hyperref[ux43aux43bux438ux43cux430ux442ux438ux447ux435ux441ux43aux438ux435-ux443ux441ux43bux43eux432ux438ux44f-ux44dux43aux441ux43fux43bux443ux430ux442ux430ux446ux438ux438]{4.3.1
Климатические условия эксплуатации 14}

\hyperref[ux442ux440ux435ux431ux43eux432ux430ux43dux438ux44f-ux43a-ux43aux432ux430ux43bux438ux444ux438ux43aux430ux446ux438ux438-ux438-ux447ux438ux441ux43bux435ux43dux43dux43eux441ux442ux438-ux43fux435ux440ux441ux43eux43dux430ux43bux430]{4.3.2
Требования к квалификации и численности персонала 14}

\hyperref[ux442ux440ux435ux431ux43eux432ux430ux43dux438ux44f-ux43a-ux441ux43eux441ux442ux430ux432ux443-ux438-ux43fux430ux440ux430ux43cux435ux442ux440ux430ux43c-ux442ux435ux445ux43dux438ux447ux435ux441ux43aux438ux445-ux441ux440ux435ux434ux441ux442ux432]{4.4
Требования к составу и параметрам технических средств 14}

\hyperref[ux442ux440ux435ux431ux43eux432ux430ux43dux438ux44f-ux43a-ux438ux43dux444ux43eux440ux43cux430ux446ux438ux43eux43dux43dux43eux439-ux438-ux43fux440ux43eux433ux440ux430ux43cux43cux43dux43eux439-ux441ux43eux432ux43cux435ux441ux442ux438ux43cux43eux441ux442ux438]{4.5
Требования к информационной и программной совместимости 15}

\hyperref[ux442ux440ux435ux431ux43eux432ux430ux43dux438ux44f-ux43a-ux438ux43dux444ux43eux440ux43cux430ux446ux438ux43eux43dux43dux44bux43c-ux441ux442ux440ux443ux43aux442ux443ux440ux430ux43c-ux438-ux43cux435ux442ux43eux434ux430ux43c-ux440ux435ux448ux435ux43dux438ux44f]{4.5.1
Требования к информационным структурам и методам решения 15}

\hyperref[ux442ux440ux435ux431ux43eux432ux430ux43dux438ux44f-ux43a-ux438ux441ux445ux43eux434ux43dux44bux43c-ux43aux43eux434ux430ux43c-ux438-ux44fux437ux44bux43aux430ux43c-ux43fux440ux43eux433ux440ux430ux43cux43cux438ux440ux43eux432ux430ux43dux438ux44f]{4.5.2
Требования к исходным кодам и языкам программирования 15}

\hyperref[ux442ux440ux435ux431ux43eux432ux430ux43dux438ux44f-ux43a-ux43fux440ux43eux433ux440ux430ux43cux43cux43dux44bux43c-ux441ux440ux435ux434ux441ux442ux432ux430ux43c-ux438ux441ux43fux43eux43bux44cux437ux443ux435ux43cux44bux43c-ux43fux440ux43eux433ux440ux430ux43cux43cux43eux439]{4.5.3
Требования к программным средствам, используемым программой 15}

\hyperref[ux442ux440ux435ux431ux43eux432ux430ux43dux438ux44f-ux43a-ux437ux430ux449ux438ux442ux435-ux438ux43dux444ux43eux440ux43cux430ux446ux438ux438]{4.5.4
Требования к защите информации 15}

\hyperref[ux442ux440ux435ux431ux43eux432ux430ux43dux438ux44f-ux43a-ux43cux430ux440ux43aux438ux440ux43eux432ux43aux435-ux438-ux443ux43fux430ux43aux43eux432ux43aux435]{4.6
Требования к маркировке и упаковке 16}

\hyperref[ux442ux440ux435ux431ux43eux432ux430ux43dux438ux44f-ux43a-ux442ux440ux430ux43dux441ux43fux43eux440ux442ux438ux440ux43eux432ux430ux43dux438ux44e-ux438-ux445ux440ux430ux43dux435ux43dux438ux44e]{4.7
Требования к транспортированию и хранению 16}

\hyperref[ux441ux43fux435ux446ux438ux430ux43bux44cux43dux44bux435-ux442ux440ux435ux431ux43eux432ux430ux43dux438ux44f]{4.8
Специальные требования 16}

\hyperref[ux442ux440ux435ux431ux43eux432ux430ux43dux438ux44f-ux43a-ux43fux440ux43eux433ux440ux430ux43cux43cux43dux43eux439-ux434ux43eux43aux443ux43cux435ux43dux442ux430ux446ux438ux438]{5
Требования к программной документации 16}

\hyperref[ux442ux435ux445ux43dux438ux43aux43e-ux44dux43aux43eux43dux43eux43cux438ux447ux435ux441ux43aux438ux435-ux43fux43eux43aux430ux437ux430ux442ux435ux43bux438]{6
Технико-экономические показатели 17}

\hyperref[ux44dux43aux43eux43dux43eux43cux438ux447ux435ux441ux43aux438ux435-ux43fux440ux435ux438ux43cux443ux449ux435ux441ux442ux432ux430-ux440ux430ux437ux440ux430ux431ux43eux442ux43aux438]{6.1
Экономические преимущества разработки 17}

\hyperref[ux441ux442ux430ux434ux438ux438-ux438-ux44dux442ux430ux43fux44b-ux440ux430ux437ux440ux430ux431ux43eux442ux43aux438]{7
Стадии и этапы разработки 18}

\hyperref[ux441ux442ux430ux434ux438ux438-ux440ux430ux437ux440ux430ux431ux43eux442ux43aux438]{7.1
Стадии разработки 18}

\hyperref[ux441ux43eux434ux435ux440ux436ux430ux43dux438ux435-ux440ux430ux431ux43eux442-ux43fux43e-ux44dux442ux430ux43fux430ux43c]{7.2
Содержание работ по этапам 19}

\hyperref[_mqq6yz6vxe9d]{8 Порядок контроля и приемки 19}

\hyperref[ux432ux438ux434ux44b-ux438ux441ux43fux44bux442ux430ux43dux438ux439]{8.1
Виды испытаний 20}

\hyperref[ux43fux440ux438ux43bux43eux436ux435ux43dux438ux435-ux431.1]{Приложение
Б.1} --
\hyperref[ux434ux438ux430ux433ux440ux430ux43cux43cux430-ux432ux430ux440ux438ux430ux43dux442ux43eux432-ux438ux441ux43fux43eux43bux44cux437ux43eux432ux430ux43dux438ux44f]{Диаграмма
вариантов использования 21}

\hyperref[ux43fux440ux438ux43bux43eux436ux435ux43dux438ux435-ux431.2]{Приложение
Б.2} --
\hyperref[ux441ux446ux435ux43dux430ux440ux438ux438-ux432ux430ux440ux438ux430ux43dux442ux43eux432-ux438ux441ux43fux43eux43bux44cux437ux43eux432ux430ux43dux438ux44f]{Сценарии
вариантов использования 22}

\hyperref[ux43fux440ux438ux43bux43eux436ux435ux43dux438ux435-ux431.3]{Приложение
Б.3}
\hyperref[ux43fux440ux438ux43bux43eux436ux435ux43dux438ux435-ux431.3]{--}
\hyperref[ux43cux430ux43aux435ux442ux44b-ux44dux43aux440ux430ux43dux43dux44bux445-ux444ux43eux440ux43c]{Макеты
экранных форм 25}

\hyperref[ux43fux440ux438ux43bux43eux436ux435ux43dux438ux435-ux431.4]{Приложение
Б.4}
\hyperref[ux43fux440ux438ux43bux43eux436ux435ux43dux438ux435-ux431.4]{--}
\hyperref[ux441ux442ux440ux443ux43aux442ux443ux440ux430-ux438-ux444ux43eux440ux43cux430ux442-ux434ux430ux43dux43dux44bux445]{Структура
и формат данных 27}

\hyperref[ux43fux440ux438ux43bux43eux436ux435ux43dux438ux435-ux431.5]{Приложение
Б.5}
\hyperref[ux43fux440ux438ux43bux43eux436ux435ux43dux438ux435-ux431.5]{--}
\hyperref[ux441ux43eux441ux442ux43eux44fux43dux438ux44f-ux444ux438ux437ux438ux447ux435ux441ux43aux43eux433ux43e-ux443ux440ux43eux432ux43dux44f-ux43cux43eux434ux435ux43bux438-ux43fux440ux438ux451ux43cux43eux43fux435ux440ux435ux434ux430ux442ux447ux438ux43aux430]{Состояния
физического уровня модели приёмопередатчика 28}

\section{Введение}\label{ux432ux432ux435ux434ux435ux43dux438ux435}

\section{1.1 Наименование
программы}\label{ux43dux430ux438ux43cux435ux43dux43eux432ux430ux43dux438ux435-ux43fux440ux43eux433ux440ux430ux43cux43cux44b}

Наименование программы: ``Программный модуль модели гидроакустического
приёмопередатчика для Ns-3''.

\section{1.2 Область
применения}\label{ux43eux431ux43bux430ux441ux442ux44c-ux43fux440ux438ux43cux435ux43dux435ux43dux438ux44f}

Модель применяется для:

\begin{itemize}
\item
  Виртуальных испытаний и исследовательского моделирования
  гидроакустических приёмопередатчиков и каналов связи;
\item
  Воспроизведения сценариев передачи сообщений между подводными
  объектами/узлами в условиях, характерных для гидроакустики;
\item
  Отработки алгоритмов верхнего уровня на базе результатов моделирования
  канала и приёмопередатчика.
\end{itemize}

При этом модуль предназначен для предоставления сетевой и физической
модели гидроакустического канала связи и обмена сообщениями между узлами
в среде ns-3. Алгоритмы верхнего уровня в состав модуля не входят и
могут разрабатываться и отрабатываться на его базе с использованием
предоставляемых средств моделирования.

Модуль ориентирован на применение инженерами и исследователями
организации, занимающейся гидроакустикой, на рабочих станциях
специалистов.

\section{2. Основание для
разработки}\label{ux43eux441ux43dux43eux432ux430ux43dux438ux435-ux434ux43bux44f-ux440ux430ux437ux440ux430ux431ux43eux442ux43aux438}

Разработка программы ведётся на основании задания на выпускную работу
бакалавра, полученного в соответствии с приказом №1203-ст от 05 сентября
2025 года «Об утверждении тем и руководителей выпускных работ
бакалавров» по кафедре ПОАС на тему «Разработка компьютерной модели
программно-аппаратного гидроакустического приёмопередатчика для его
виртуальных испытаний» и согласовано с ФГУП ВНИИА им. Духова НПЦ
``Гидросвязь''.

\section{3. Назначение
разработки}\label{ux43dux430ux437ux43dux430ux447ux435ux43dux438ux435-ux440ux430ux437ux440ux430ux431ux43eux442ux43aux438}

Цель разработки: сократить расходы и время на разработку систем
гидроакустической связи за счёт разработки компьютерной модели
программно-аппаратного гидроакустического приёмопередатчика для
виртуальных испытаний систем гидроакустической связи и отработки
алгоритмов функционирования мультиагентных систем.

Программный модуль предназначен для эксплуатации инженерами и
программистами в ФГУП ВНИИА им. Духова НПЦ ``Гидросвязь''.

\section{4 Требования к
программе}\label{ux442ux440ux435ux431ux43eux432ux430ux43dux438ux44f-ux43a-ux43fux440ux43eux433ux440ux430ux43cux43cux435}

\section{4.1 Требования к функциональным
характеристикам}\label{ux442ux440ux435ux431ux43eux432ux430ux43dux438ux44f-ux43a-ux444ux443ux43dux43aux446ux438ux43eux43dux430ux43bux44cux43dux44bux43c-ux445ux430ux440ux430ux43aux442ux435ux440ux438ux441ux442ux438ux43aux430ux43c}

4.1.1 Состав программного модуля

Программный модуль должен включать:

\begin{itemize}
\item
  Модель гидроакустического приёмопередатчика;
\item
  Модель гидроакустического канала связи;
\item
  Набор базовых компонентов, необходимых для реализации других моделей
  гидроакустических приёмопередатчиков;
\item
  Набор базовых компонентов, необходимых для реализации других моделей
  гидроакустических каналов связи.
\end{itemize}

Программный модуль должен содержать не только конкретные реализации
моделей, но и предлагать компоненты и архитектурные подходы, необходимые
для реализации других моделей.

\section{4.1.2 Требования к архитектуре модели программно-аппаратного
приёмопередатчика}\label{ux442ux440ux435ux431ux43eux432ux430ux43dux438ux44f-ux43a-ux430ux440ux445ux438ux442ux435ux43aux442ux443ux440ux435-ux43cux43eux434ux435ux43bux438-ux43fux440ux43eux433ux440ux430ux43cux43cux43dux43e-ux430ux43fux43fux430ux440ux430ux442ux43dux43eux433ux43e-ux43fux440ux438ux451ux43cux43eux43fux435ux440ux435ux434ux430ux442ux447ux438ux43aux430}

\section{4.1.2.1 Общие
требования}\label{ux43eux431ux449ux438ux435-ux442ux440ux435ux431ux43eux432ux430ux43dux438ux44f}

Архитектура модели приёмопередатчика должна обеспечивать реализацию
физического и канального уровней устройства связи по модели OSI.

Архитектура модели приёмопередатчика должна обеспечивать совместимость с
моделями каналов гидроакустической связи.

Архитектура модели приёмопередатчика должна обеспечивать возможность
реализации конкретных моделей приёмопередатчиков без необходимости
изменения интерфейсов базовых классов.

\section{4.1.2.2 Независимость
уровней}\label{ux43dux435ux437ux430ux432ux438ux441ux438ux43cux43eux441ux442ux44c-ux443ux440ux43eux432ux43dux435ux439}

Физический и канальный уровень приёмопередатчика должны быть разделены.

Взаимодействие между уровнями должно быть организовано через чётко
определённые интерфейсы и регламентированные механизмы обратного вызова.

Архитектура физического уровня должна быть независима от архитектуры
канального уровня.

\section{4.1.2.3 Функции физического
уровня}\label{ux444ux443ux43dux43aux446ux438ux438-ux444ux438ux437ux438ux447ux435ux441ux43aux43eux433ux43e-ux443ux440ux43eux432ux43dux44f}

Физический уровень приёмопередатчика должен обеспечивать:

\begin{itemize}
\item
  Формирование передаваемого сигнала на основе входящего информационного
  пакета;
\item
  \hl{Сигнал на физическом уровне описывается спектральной плотностью
  мощности и его продолжительностью;}
\item
  Передачу сигнала в модель гидроакустического канала связи;
\item
  Приём сигнала из модели гидроакустического канала связи;
\item
  Обработку принятого сигнала и извлечение полезной нагрузки (полезной
  информации) из принятого сигнала.
\end{itemize}

\section{4.1.2.4 Модели
интерференции}\label{ux43cux43eux434ux435ux43bux438-ux438ux43dux442ux435ux440ux444ux435ux440ux435ux43dux446ux438ux438}

Архитектура приёмопередатчика должна поддерживать работу с различными
моделями интерференции.

Модель интерференции должна рассчитывать значение отношения сигнал/шум
(ОСШ) для каждого входящего сигнала, который потенциально может быть
принят и декодирован.

Решение о том, будет ли корректно принят и декодирован входящий сигнал,
принимается в конкретных реализациях моделей приёмопередатчиков, в том
числе на основании рассчитанных моделями интерференции характеристик
качества приёма сигнала.

\section{4.1.3 Требования к архитектуре модели гидроакустического канала
связи}\label{ux442ux440ux435ux431ux43eux432ux430ux43dux438ux44f-ux43a-ux430ux440ux445ux438ux442ux435ux43aux442ux443ux440ux435-ux43cux43eux434ux435ux43bux438-ux433ux438ux434ux440ux43eux430ux43aux443ux441ux442ux438ux447ux435ux441ux43aux43eux433ux43e-ux43aux430ux43dux430ux43bux430-ux441ux432ux44fux437ux438}

Архитектура модели гидроакустического канала связи должна обеспечивать
взаимодействие с физическим уровнем модели приёмопередатчика, включая
передачу сигналов от физического уровня в модель среды распространения и
обратно.

Архитектура канала должна предусматривать возможность создания
конкретных реализаций моделей каналов связи.

Расчёт таких эффектов, как задержка, затухание, многолучёвость, эффект
Доплера, должен осуществляться через регламентированные интерфейсы.
Модели каналов связи для расчёта свойств распространения сигналов должны
использовать эти интерфейсы.\hl{}

\hl{}

\section{4.1.4 Требования к реализации модели
приёмопередатчика}\label{ux442ux440ux435ux431ux43eux432ux430ux43dux438ux44f-ux43a-ux440ux435ux430ux43bux438ux437ux430ux446ux438ux438-ux43cux43eux434ux435ux43bux438-ux43fux440ux438ux451ux43cux43eux43fux435ux440ux435ux434ux430ux442ux447ux438ux43aux430}

\hl{}

\section{4.1.4.1 Общие
положения}\label{ux43eux431ux449ux438ux435-ux43fux43eux43bux43eux436ux435ux43dux438ux44f}

Модель приёмопередатчика должна быть основана на архитектуре
программно-аппаратного гидроакустического приёмопередатчика, требования
к которой были описаны ранее.

Модель приёмопередатчика должна реализовывать физический и канальный
уровень модели OSI.

Модель приёмопередатчика поддерживает только один формат передаваемого и
принимаемого сигнала. Сигналы другого типа приёмопередатчик принимать не
должен.

\section{4.1.4.2 Состояния физического
уровня}\label{ux441ux43eux441ux442ux43eux44fux43dux438ux44f-ux444ux438ux437ux438ux447ux435ux441ux43aux43eux433ux43e-ux443ux440ux43eux432ux43dux44f}

Физический уровень модели приёмопередатчика должен реализовывать
конечный автомат со следующими состояниями:

\begin{itemize}
\item
  OFF -- устройство отключено (передача и приём не выполняются);
\item
  IDLE -- устройство находится в режиме ожидания и осуществляет
  прослушивание канала, готово к отправке и приёму сигналов;
\item
  RX -- устройство находится в состоянии приёма и обработки входящего
  сигнала;
\item
  TX -- устройство находится в состоянии передачи сигнала.
\end{itemize}

Физический уровень модели должен находиться только в одном состоянии в
каждый момент модельного времени.

Базовая реализация приёмопередатчика должна обеспечивать режим
half-duplex: одновременные передача и приём не допускаются.

Подробное описание поведения конечного автомата физического уровня
приведено на диаграмме состояний, представленной в приложении Б.5.

\section{4.1.4.3 Функции канального
уровня}\label{ux444ux443ux43dux43aux446ux438ux438-ux43aux430ux43dux430ux43bux44cux43dux43eux433ux43e-ux443ux440ux43eux432ux43dux44f}

Канальный уровень модели приёмопередатчика должен обеспечивать
буферизацию данных на передачу в виде очереди пакетов. Пользователь
должен иметь возможность задать максимальный размер очереди
(определяется количеством пакетов). При достижении максимального размера
очереди, новые пакеты должны быть отброшены.

Приёмопередатчик должен иметь MAC-адрес (уникальный идентификатор
канального уровня).

\hl{MAC-адрес состоит из 8 бит.}

Канальный уровень модели приёмопередатчика должен добавлять к
передаваемым данным заголовок, содержащий:

\begin{itemize}
\item
  MAC-адрес отправителя;
\item
  MAC-адрес получателя.
\end{itemize}

\section{4.1.4.4 Используемая модель
интерференции}\label{ux438ux441ux43fux43eux43bux44cux437ux443ux435ux43cux430ux44f-ux43cux43eux434ux435ux43bux44c-ux438ux43dux442ux435ux440ux444ux435ux440ux435ux43dux446ux438ux438}

\hl{Модель интерференции должна вычислять отношение сигнал к шуму по
следующим формулам} (см.~формулы~1-4):\hl{}

\begin{longtable}[]{@{}
  >{\raggedright\arraybackslash}p{(\columnwidth - 4\tabcolsep) * \real{0.0523}}
  >{\raggedright\arraybackslash}p{(\columnwidth - 4\tabcolsep) * \real{0.8990}}
  >{\raggedright\arraybackslash}p{(\columnwidth - 4\tabcolsep) * \real{0.0487}}@{}}
\toprule\noalign{}
\multicolumn{3}{@{}>{\raggedright\arraybackslash}p{(\columnwidth - 4\tabcolsep) * \real{1.0000} + 4\tabcolsep}@{}}{%
\begin{minipage}[b]{\linewidth}\raggedright
\end{minipage}} \\
\midrule\noalign{}
\endhead
\bottomrule\noalign{}
\endlastfoot
& \(P_{s} = \int_{f_{1}}^{f_{2}}{S_{rx}(f){df}}\) & (1) \\
\multicolumn{3}{@{}>{\raggedright\arraybackslash}p{(\columnwidth - 4\tabcolsep) * \real{1.0000} + 4\tabcolsep}@{}}{%
} \\
где~ &
\multicolumn{2}{>{\raggedright\arraybackslash}p{(\columnwidth - 4\tabcolsep) * \real{0.9477} + 2\tabcolsep}@{}}{%
\(P_{s}\) -- мощность полезного сигнала в рабочей полосе, Па²;} \\
&
\multicolumn{2}{>{\raggedright\arraybackslash}p{(\columnwidth - 4\tabcolsep) * \real{0.9477} + 2\tabcolsep}@{}}{%
\(f_{1}\) -- нижняя граница рабочей полосы, Гц;} \\
&
\multicolumn{2}{>{\raggedright\arraybackslash}p{(\columnwidth - 4\tabcolsep) * \real{0.9477} + 2\tabcolsep}@{}}{%
\(f_{2}\) -- верхняя граница рабочей полосы, Гц;} \\
&
\multicolumn{2}{>{\raggedright\arraybackslash}p{(\columnwidth - 4\tabcolsep) * \real{0.9477} + 2\tabcolsep}@{}}{%
\(S_{rx}\) -- спектральная плотность мощности полезного сигнала на входе
приёмника, Па²/Гц.} \\
\end{longtable}

\begin{longtable}[]{@{}
  >{\raggedright\arraybackslash}p{(\columnwidth - 4\tabcolsep) * \real{0.0523}}
  >{\raggedright\arraybackslash}p{(\columnwidth - 4\tabcolsep) * \real{0.8990}}
  >{\raggedright\arraybackslash}p{(\columnwidth - 4\tabcolsep) * \real{0.0487}}@{}}
\toprule\noalign{}
\multicolumn{3}{@{}>{\raggedright\arraybackslash}p{(\columnwidth - 4\tabcolsep) * \real{1.0000} + 4\tabcolsep}@{}}{%
\begin{minipage}[b]{\linewidth}\raggedright
\end{minipage}} \\
\midrule\noalign{}
\endhead
\bottomrule\noalign{}
\endlastfoot
& \(P_{n} = \int_{f_{1}}^{f_{2}}{S_{n}(f){df}}\) & (2) \\
& \(P_{i} = \int_{f_{1}}^{f_{2}}{S_{i}(f){df}}\) & (3) \\
\multicolumn{3}{@{}>{\raggedright\arraybackslash}p{(\columnwidth - 4\tabcolsep) * \real{1.0000} + 4\tabcolsep}@{}}{%
} \\
где~ &
\multicolumn{2}{>{\raggedright\arraybackslash}p{(\columnwidth - 4\tabcolsep) * \real{0.9477} + 2\tabcolsep}@{}}{%
\(P_{n}\) -- мощность шума в рабочей полосе, Па²;} \\
&
\multicolumn{2}{>{\raggedright\arraybackslash}p{(\columnwidth - 4\tabcolsep) * \real{0.9477} + 2\tabcolsep}@{}}{%
\(S_{n}\) -- спектральная плотность мощности шума, Па²/Гц;} \\
&
\multicolumn{2}{>{\raggedright\arraybackslash}p{(\columnwidth - 4\tabcolsep) * \real{0.9477} + 2\tabcolsep}@{}}{%
\(P_{i}\) -- мощность помехи (другого сигнала) в рабочей полосе,
Па²;} \\
&
\multicolumn{2}{>{\raggedright\arraybackslash}p{(\columnwidth - 4\tabcolsep) * \real{0.9477} + 2\tabcolsep}@{}}{%
\(S_{i}\) -- спектральная плотность мощности помехи на входе приёмника,
Па²/Гц.} \\
\end{longtable}

\hl{}

\begin{longtable}[]{@{}
  >{\raggedright\arraybackslash}p{(\columnwidth - 4\tabcolsep) * \real{0.0523}}
  >{\raggedright\arraybackslash}p{(\columnwidth - 4\tabcolsep) * \real{0.8990}}
  >{\raggedright\arraybackslash}p{(\columnwidth - 4\tabcolsep) * \real{0.0487}}@{}}
\toprule\noalign{}
\multicolumn{3}{@{}>{\raggedright\arraybackslash}p{(\columnwidth - 4\tabcolsep) * \real{1.0000} + 4\tabcolsep}@{}}{%
\begin{minipage}[b]{\linewidth}\raggedright
\end{minipage}} \\
\midrule\noalign{}
\endhead
\bottomrule\noalign{}
\endlastfoot
& \(SINR = \frac{P_{s}}{\sum_{k = 1}^{N_{i}}{(P_{i,k})}\  + \ P_{n}}\) &
(4) \\
\multicolumn{3}{@{}>{\raggedright\arraybackslash}p{(\columnwidth - 4\tabcolsep) * \real{1.0000} + 4\tabcolsep}@{}}{%
} \\
где~ &
\multicolumn{2}{>{\raggedright\arraybackslash}p{(\columnwidth - 4\tabcolsep) * \real{0.9477} + 2\tabcolsep}@{}}{%
\(SINR\) -- отношение мощности полезного сигнала к суммарной мощности
шума и помех;} \\
&
\multicolumn{2}{>{\raggedright\arraybackslash}p{(\columnwidth - 4\tabcolsep) * \real{0.9477} + 2\tabcolsep}@{}}{%
\(P_{s}\) -- мощность полезного сигнала в рабочей полосе, Па²;} \\
&
\multicolumn{2}{>{\raggedright\arraybackslash}p{(\columnwidth - 4\tabcolsep) * \real{0.9477} + 2\tabcolsep}@{}}{%
\(P_{n}\) -- мощность шума в рабочей полосе, Па²;} \\
&
\multicolumn{2}{>{\raggedright\arraybackslash}p{(\columnwidth - 4\tabcolsep) * \real{0.9477} + 2\tabcolsep}@{}}{%
\(N_{i}\) -- число мешающих передач (других сигналов);} \\
&
\multicolumn{2}{>{\raggedright\arraybackslash}p{(\columnwidth - 4\tabcolsep) * \real{0.9477} + 2\tabcolsep}@{}}{%
\(P_{i,k}\) -- мощность k-й помехи (другого сигнала) в рабочей полосе,
Па².} \\
\end{longtable}

\section{4.1.4.5 Решение об успешном приёме и декодировании
сигнала}\label{ux440ux435ux448ux435ux43dux438ux435-ux43eux431-ux443ux441ux43fux435ux448ux43dux43eux43c-ux43fux440ux438ux451ux43cux435-ux438-ux434ux435ux43aux43eux434ux438ux440ux43eux432ux430ux43dux438ux438-ux441ux438ux433ux43dux430ux43bux430}

Модель приёмопередатчика считает сигнал успешно принятым и
декодированным, если тип сигнала соответствует типу, который может
декодировать этот приёмопередатчик, и значение ОСШ для этого сигнала
выше порогового значения, заданного пользователем (см.~формулу~5):

\begin{longtable}[]{@{}
  >{\raggedright\arraybackslash}p{(\columnwidth - 4\tabcolsep) * \real{0.0523}}
  >{\raggedright\arraybackslash}p{(\columnwidth - 4\tabcolsep) * \real{0.8990}}
  >{\raggedright\arraybackslash}p{(\columnwidth - 4\tabcolsep) * \real{0.0487}}@{}}
\toprule\noalign{}
\multicolumn{3}{@{}>{\raggedright\arraybackslash}p{(\columnwidth - 4\tabcolsep) * \real{1.0000} + 4\tabcolsep}@{}}{%
\begin{minipage}[b]{\linewidth}\raggedright
\end{minipage}} \\
\midrule\noalign{}
\endhead
\bottomrule\noalign{}
\endlastfoot
& \(SINR \geq \gamma_{th}\) & (5) \\
\multicolumn{3}{@{}>{\raggedright\arraybackslash}p{(\columnwidth - 4\tabcolsep) * \real{1.0000} + 4\tabcolsep}@{}}{%
} \\
где~ &
\multicolumn{2}{>{\raggedright\arraybackslash}p{(\columnwidth - 4\tabcolsep) * \real{0.9477} + 2\tabcolsep}@{}}{%
\(SINR\) -- отношение мощности полезного сигнала к суммарной мощности
шума и помех;} \\
&
\multicolumn{2}{>{\raggedright\arraybackslash}p{(\columnwidth - 4\tabcolsep) * \real{0.9477} + 2\tabcolsep}@{}}{%
\(\gamma_{th}\) -- порог успешного приёма.} \\
\end{longtable}

\hl{}

\section{4.1.4.6
Трассировка}\label{ux442ux440ux430ux441ux441ux438ux440ux43eux432ux43aux430}

Модель приёмопередатчика должна поддерживать вывод результатов
трассировки в формате pcap. Файл должен содержать информацию о принятых
и отправленных пакетах.

\section{4.1.5 Требования к модели гидроакустического канала
связи}\label{ux442ux440ux435ux431ux43eux432ux430ux43dux438ux44f-ux43a-ux43cux43eux434ux435ux43bux438-ux433ux438ux434ux440ux43eux430ux43aux443ux441ux442ux438ux447ux435ux441ux43aux43eux433ux43e-ux43aux430ux43dux430ux43bux430-ux441ux432ux44fux437ux438}

\section{4.1.5.1 Расчёт характеристик
распространения}\label{ux440ux430ux441ux447ux451ux442-ux445ux430ux440ux430ux43aux442ux435ux440ux438ux441ux442ux438ux43a-ux440ux430ux441ux43fux440ux43eux441ux442ux440ux430ux43dux435ux43dux438ux44f}

Базовая модель канала связи должна обеспечивать:

\begin{itemize}
\item
  расчёт времени распространения сигнала от точки излучения до точки
  приёма;
\item
  расчёт затухания (потерь) сигнала при прохождении из точки излучения в
  точку приёма.
\end{itemize}

\section{4.1.5.2 Расчёт времени
распространения}\label{ux440ux430ux441ux447ux451ux442-ux432ux440ux435ux43cux435ux43dux438-ux440ux430ux441ux43fux440ux43eux441ux442ux440ux430ux43dux435ux43dux438ux44f}

Базовая модель канала должна осуществлять расчёт времени распространения
сигнала на основании 2 заданных параметров:

\begin{itemize}
\item
  Координата излучателя в трёхмерном пространстве;
\item
  Координата приёмника в трёхмерном пространстве.
\end{itemize}

Для расчёта времени распространения считать:

\begin{itemize}
\item
  Скорость распространения звука в воде постоянная;
\item
  Сигнал однолучевой;
\item
  Расстояние между узлами определяется как евклидово расстояние между их
  координатами (см. формулу 6).
\end{itemize}

\begin{longtable}[]{@{}
  >{\raggedright\arraybackslash}p{(\columnwidth - 4\tabcolsep) * \real{0.0523}}
  >{\raggedright\arraybackslash}p{(\columnwidth - 4\tabcolsep) * \real{0.8990}}
  >{\raggedright\arraybackslash}p{(\columnwidth - 4\tabcolsep) * \real{0.0487}}@{}}
\toprule\noalign{}
\multicolumn{3}{@{}>{\raggedright\arraybackslash}p{(\columnwidth - 4\tabcolsep) * \real{1.0000} + 4\tabcolsep}@{}}{%
\begin{minipage}[b]{\linewidth}\raggedright
\end{minipage}} \\
\midrule\noalign{}
\endhead
\bottomrule\noalign{}
\endlastfoot
& \(d = \left\| r_{rx} - r_{tx} \right\|\) & (6) \\
\multicolumn{3}{@{}>{\raggedright\arraybackslash}p{(\columnwidth - 4\tabcolsep) * \real{1.0000} + 4\tabcolsep}@{}}{%
} \\
где~ &
\multicolumn{2}{>{\raggedright\arraybackslash}p{(\columnwidth - 4\tabcolsep) * \real{0.9477} + 2\tabcolsep}@{}}{%
\(d\) -- расстояние между передатчиком и приёмником, м;} \\
&
\multicolumn{2}{>{\raggedright\arraybackslash}p{(\columnwidth - 4\tabcolsep) * \real{0.9477} + 2\tabcolsep}@{}}{%
\(r_{tx}\) -- радиус-вектор передатчика;} \\
&
\multicolumn{2}{>{\raggedright\arraybackslash}p{(\columnwidth - 4\tabcolsep) * \real{0.9477} + 2\tabcolsep}@{}}{%
\(r_{rx}\) -- радиус-вектор приёмника.} \\
\end{longtable}

Время распространения определяется по формуле 7:

\begin{longtable}[]{@{}
  >{\raggedright\arraybackslash}p{(\columnwidth - 4\tabcolsep) * \real{0.0523}}
  >{\raggedright\arraybackslash}p{(\columnwidth - 4\tabcolsep) * \real{0.8990}}
  >{\raggedright\arraybackslash}p{(\columnwidth - 4\tabcolsep) * \real{0.0487}}@{}}
\toprule\noalign{}
\multicolumn{3}{@{}>{\raggedright\arraybackslash}p{(\columnwidth - 4\tabcolsep) * \real{1.0000} + 4\tabcolsep}@{}}{%
\begin{minipage}[b]{\linewidth}\raggedright
\end{minipage}} \\
\midrule\noalign{}
\endhead
\bottomrule\noalign{}
\endlastfoot
& \(\tau = \frac{d}{c}\) & (7) \\
\multicolumn{3}{@{}>{\raggedright\arraybackslash}p{(\columnwidth - 4\tabcolsep) * \real{1.0000} + 4\tabcolsep}@{}}{%
} \\
где~ &
\multicolumn{2}{>{\raggedright\arraybackslash}p{(\columnwidth - 4\tabcolsep) * \real{0.9477} + 2\tabcolsep}@{}}{%
\(\tau\) -- задержка распространения сигнала, с.} \\
&
\multicolumn{2}{>{\raggedright\arraybackslash}p{(\columnwidth - 4\tabcolsep) * \real{0.9477} + 2\tabcolsep}@{}}{%
\(d\) -- расстояние между передатчиком и приёмником, м.} \\
&
\multicolumn{2}{>{\raggedright\arraybackslash}p{(\columnwidth - 4\tabcolsep) * \real{0.9477} + 2\tabcolsep}@{}}{%
\(c\) -- скорость звука в воде, м/с.} \\
\end{longtable}

\section{4.1.5.3 Расчёт затухания (потерь)
сигнала}\label{ux440ux430ux441ux447ux451ux442-ux437ux430ux442ux443ux445ux430ux43dux438ux44f-ux43fux43eux442ux435ux440ux44c-ux441ux438ux433ux43dux430ux43bux430}

Базовая модель канала должна осуществлять расчёт затухания сигнала на
основании 3 заданных параметров:

\begin{itemize}
\item
  координата излучателя в трёхмерном пространстве;
\item
  координата приёмника в трёхмерном пространстве;
\item
  мощность сигнала на выходе излучателя.
\end{itemize}

Модель потерь состоит из нескольких составляющих (см. формулу 8):

\begin{longtable}[]{@{}
  >{\raggedright\arraybackslash}p{(\columnwidth - 4\tabcolsep) * \real{0.0523}}
  >{\raggedright\arraybackslash}p{(\columnwidth - 4\tabcolsep) * \real{0.8990}}
  >{\raggedright\arraybackslash}p{(\columnwidth - 4\tabcolsep) * \real{0.0487}}@{}}
\toprule\noalign{}
\multicolumn{3}{@{}>{\raggedright\arraybackslash}p{(\columnwidth - 4\tabcolsep) * \real{1.0000} + 4\tabcolsep}@{}}{%
\begin{minipage}[b]{\linewidth}\raggedright
\end{minipage}} \\
\midrule\noalign{}
\endhead
\bottomrule\noalign{}
\endlastfoot
& \(TL(f,d) = {TL}_{geom}(d) + {TL}_{abs}(f,d)\) & (8) \\
\multicolumn{3}{@{}>{\raggedright\arraybackslash}p{(\columnwidth - 4\tabcolsep) * \real{1.0000} + 4\tabcolsep}@{}}{%
} \\
где~ &
\multicolumn{2}{>{\raggedright\arraybackslash}p{(\columnwidth - 4\tabcolsep) * \real{0.9477} + 2\tabcolsep}@{}}{%
\(TL\) (Transmission loss) -- суммарные потери сигнала на частоте \(f\)
при расстоянии \(d\), дБ;} \\
&
\multicolumn{2}{>{\raggedright\arraybackslash}p{(\columnwidth - 4\tabcolsep) * \real{0.9477} + 2\tabcolsep}@{}}{%
\({TL}_{geom}\) -- потери на геометрическое расхождение, дБ;} \\
&
\multicolumn{2}{>{\raggedright\arraybackslash}p{(\columnwidth - 4\tabcolsep) * \real{0.9477} + 2\tabcolsep}@{}}{%
\({TL}_{abs}\) -- потери на абсорбцию, дБ;} \\
&
\multicolumn{2}{>{\raggedright\arraybackslash}p{(\columnwidth - 4\tabcolsep) * \real{0.9477} + 2\tabcolsep}@{}}{%
\(f\) -- частота сигнала, Гц;} \\
&
\multicolumn{2}{>{\raggedright\arraybackslash}p{(\columnwidth - 4\tabcolsep) * \real{0.9477} + 2\tabcolsep}@{}}{%
\(d\) -- расстояние, пройденное сигналом, м.} \\
\end{longtable}

Потери на геометрическое расхождение вычисляются по формуле 9:

\begin{longtable}[]{@{}
  >{\raggedright\arraybackslash}p{(\columnwidth - 4\tabcolsep) * \real{0.0523}}
  >{\raggedright\arraybackslash}p{(\columnwidth - 4\tabcolsep) * \real{0.8987}}
  >{\raggedright\arraybackslash}p{(\columnwidth - 4\tabcolsep) * \real{0.0490}}@{}}
\toprule\noalign{}
\multicolumn{3}{@{}>{\raggedright\arraybackslash}p{(\columnwidth - 4\tabcolsep) * \real{1.0000} + 4\tabcolsep}@{}}{%
\begin{minipage}[b]{\linewidth}\raggedright
\end{minipage}} \\
\midrule\noalign{}
\endhead
\bottomrule\noalign{}
\endlastfoot
& \({TL}_{geom}(d) = 10n\log_{10}d\)~ & (9) \\
\multicolumn{3}{@{}>{\raggedright\arraybackslash}p{(\columnwidth - 4\tabcolsep) * \real{1.0000} + 4\tabcolsep}@{}}{%
} \\
где~ &
\multicolumn{2}{>{\raggedright\arraybackslash}p{(\columnwidth - 4\tabcolsep) * \real{0.9477} + 2\tabcolsep}@{}}{%
\({TL}_{geom}\) -- потери на геометрическое расхождение, дБ;} \\
&
\multicolumn{2}{>{\raggedright\arraybackslash}p{(\columnwidth - 4\tabcolsep) * \real{0.9477} + 2\tabcolsep}@{}}{%
\(n\) -- коэффициент расхождения;} \\
&
\multicolumn{2}{>{\raggedright\arraybackslash}p{(\columnwidth - 4\tabcolsep) * \real{0.9477} + 2\tabcolsep}@{}}{%
\(d\) -- расстояние, пройденное сигналом, м.} \\
\end{longtable}

\hl{Коэффициент n является настраиваемым и находится в диапазоне
n~∈~{[}1;2{]}.}

Потери на абсорбцию вычисляются по формуле 10:\hl{}

\begin{longtable}[]{@{}
  >{\raggedright\arraybackslash}p{(\columnwidth - 4\tabcolsep) * \real{0.0521}}
  >{\raggedright\arraybackslash}p{(\columnwidth - 4\tabcolsep) * \real{0.8799}}
  >{\raggedright\arraybackslash}p{(\columnwidth - 4\tabcolsep) * \real{0.0680}}@{}}
\toprule\noalign{}
\multicolumn{3}{@{}>{\raggedright\arraybackslash}p{(\columnwidth - 4\tabcolsep) * \real{1.0000} + 4\tabcolsep}@{}}{%
\begin{minipage}[b]{\linewidth}\raggedright
\end{minipage}} \\
\midrule\noalign{}
\endhead
\bottomrule\noalign{}
\endlastfoot
& \({TL}_{abs}(f,\ d) = \alpha(f)*d*10^{3}\) & (10) \\
\multicolumn{3}{@{}>{\raggedright\arraybackslash}p{(\columnwidth - 4\tabcolsep) * \real{1.0000} + 4\tabcolsep}@{}}{%
} \\
где~ &
\multicolumn{2}{>{\raggedright\arraybackslash}p{(\columnwidth - 4\tabcolsep) * \real{0.9479} + 2\tabcolsep}@{}}{%
\({TL}_{abs}\) -- потери на абсорбцию, дБ;} \\
&
\multicolumn{2}{>{\raggedright\arraybackslash}p{(\columnwidth - 4\tabcolsep) * \real{0.9479} + 2\tabcolsep}@{}}{%
\(\alpha\) -- коэффициент поглощения, дБ/км;} \\
&
\multicolumn{2}{>{\raggedright\arraybackslash}p{(\columnwidth - 4\tabcolsep) * \real{0.9479} + 2\tabcolsep}@{}}{%
\(d\) -- расстояние, пройденное сигналом, м.} \\
\end{longtable}

\hl{Для вычисления коэффициента поглощения используются формулы Торпа}
(см.~формулу~11) и Франсуа-Гаррисона\hl{} (см.~формулу~12):\hl{}

\begin{longtable}[]{@{}
  >{\raggedright\arraybackslash}p{(\columnwidth - 4\tabcolsep) * \real{0.0523}}
  >{\raggedright\arraybackslash}p{(\columnwidth - 4\tabcolsep) * \real{0.8826}}
  >{\raggedright\arraybackslash}p{(\columnwidth - 4\tabcolsep) * \real{0.0651}}@{}}
\toprule\noalign{}
\multicolumn{3}{@{}>{\raggedright\arraybackslash}p{(\columnwidth - 4\tabcolsep) * \real{1.0000} + 4\tabcolsep}@{}}{%
\begin{minipage}[b]{\linewidth}\raggedright
\end{minipage}} \\
\midrule\noalign{}
\endhead
\bottomrule\noalign{}
\endlastfoot
&
\(\alpha(f) \approx \frac{0.11\ f^{2}}{1\  + \ f^{2}} + \frac{40f^{2}}{4100 + f^{2}} + \frac{0,03f^{2}}{36 + f^{2}}\ \)
& (11) \\
\multicolumn{3}{@{}>{\raggedright\arraybackslash}p{(\columnwidth - 4\tabcolsep) * \real{1.0000} + 4\tabcolsep}@{}}{%
} \\
где~ &
\multicolumn{2}{>{\raggedright\arraybackslash}p{(\columnwidth - 4\tabcolsep) * \real{0.9477} + 2\tabcolsep}@{}}{%
\(\alpha\) -- коэффициент поглощения, дБ/км;} \\
&
\multicolumn{2}{>{\raggedright\arraybackslash}p{(\columnwidth - 4\tabcolsep) * \real{0.9477} + 2\tabcolsep}@{}}{%
\(f\) -- частота звука, кГц.} \\
\end{longtable}

\begin{longtable}[]{@{}
  >{\raggedright\arraybackslash}p{(\columnwidth - 4\tabcolsep) * \real{0.0661}}
  >{\raggedright\arraybackslash}p{(\columnwidth - 4\tabcolsep) * \real{0.8684}}
  >{\raggedright\arraybackslash}p{(\columnwidth - 4\tabcolsep) * \real{0.0654}}@{}}
\toprule\noalign{}
\multicolumn{3}{@{}>{\raggedright\arraybackslash}p{(\columnwidth - 4\tabcolsep) * \real{1.0000} + 4\tabcolsep}@{}}{%
\begin{minipage}[b]{\linewidth}\raggedright
\end{minipage}} \\
\midrule\noalign{}
\endhead
\bottomrule\noalign{}
\endlastfoot
&
\(\alpha(f) \approx \frac{A_{1}P_{1}f_{1}f^{2}}{f_{1}^{2} + f^{2}} + \frac{A_{2}P_{2}f_{2}f^{2}}{f_{2}^{2} + f^{2}} + A_{3}P_{3}f^{2}\)
& (12) \\
& \(\begin{matrix}
A_{1} = \frac{8.86}{c} \cdot {10}^{0.78pH - 5}
\end{matrix}\), \(P_{1} = 1\)

\(\begin{matrix}
f_{1} = 2.8\sqrt[]{\frac{S}{35}} \cdot 10^{4\  - \ \frac{1245}{273 + T}}
\end{matrix}\)

\(\begin{matrix}
A_{2} = 21.44\frac{S}{c}(1 + 0.025T)
\end{matrix}\)

\(\begin{matrix}
P_{2} = 1 - 1.37 \cdot 10^{- 4}D + 6.2 \cdot 10^{- 9}D^{2}
\end{matrix}\)

\(\begin{matrix}
f_{2} = \frac{8.17 \cdot 10^{8 - \frac{1990}{273 + T}}}{1 + 0.0018(S - 35)}
\end{matrix}\)

\(\begin{matrix}
A_{3} = \left\{ \begin{matrix}
4.937 \cdot 10^{- 4} - 2.59 \cdot 10^{- 5}T + 9.11 \cdot 10^{- 7}T^{2} - 1.50 \cdot 10^{- 8}T^{3}, & T \geq 20, \\
3.964 \cdot 10^{- 4} - 1.146 \cdot 10^{- 5}T + 1.45 \cdot 10^{- 7}T^{2} - 6.5 \cdot 10^{- 10}T^{3}, & T < 20.
\end{matrix} \right.\ 
\end{matrix}\)

\(\begin{matrix}
P_{3} = 1 - 3.83 \cdot 10^{- 5}D + 4.9 \cdot 10^{- 10}D^{2}
\end{matrix}\)

\(\begin{matrix}
c = 1412 + 3.21T + 1.19S + 0.0167D
\end{matrix}\) & \\
\multicolumn{3}{@{}>{\raggedright\arraybackslash}p{(\columnwidth - 4\tabcolsep) * \real{1.0000} + 4\tabcolsep}@{}}{%
} \\
где~ &
\multicolumn{2}{>{\raggedright\arraybackslash}p{(\columnwidth - 4\tabcolsep) * \real{0.9339} + 2\tabcolsep}@{}}{%
\(\alpha\) -- коэффициент поглощения, дБ/км;} \\
&
\multicolumn{2}{>{\raggedright\arraybackslash}p{(\columnwidth - 4\tabcolsep) * \real{0.9339} + 2\tabcolsep}@{}}{%
\(f\) -- частота звука, кГц;} \\
&
\multicolumn{2}{>{\raggedright\arraybackslash}p{(\columnwidth - 4\tabcolsep) * \real{0.9339} + 2\tabcolsep}@{}}{%
\(A_{1}\) -- амплитудный коэффициент вклада борной кислоты;} \\
&
\multicolumn{2}{>{\raggedright\arraybackslash}p{(\columnwidth - 4\tabcolsep) * \real{0.9339} + 2\tabcolsep}@{}}{%
\(A_{2}\) -- амплитудный коэффициент вклада сульфата магния;} \\
&
\multicolumn{2}{>{\raggedright\arraybackslash}p{(\columnwidth - 4\tabcolsep) * \real{0.9339} + 2\tabcolsep}@{}}{%
\(A_{3}\) -- амплитудный коэффициент вклада чистой воды;} \\
&
\multicolumn{2}{>{\raggedright\arraybackslash}p{(\columnwidth - 4\tabcolsep) * \real{0.9339} + 2\tabcolsep}@{}}{%
\(P_{1},\ P_{2},{\ P}_{3}\) -- поправочные множители на
давление/глубину;} \\
&
\multicolumn{2}{>{\raggedright\arraybackslash}p{(\columnwidth - 4\tabcolsep) * \real{0.9339} + 2\tabcolsep}@{}}{%
\(f_{1}\) -- релаксационная частота борной кислоты, кГц;} \\
&
\multicolumn{2}{>{\raggedright\arraybackslash}p{(\columnwidth - 4\tabcolsep) * \real{0.9339} + 2\tabcolsep}@{}}{%
\(f_{2}\) -- релаксационная частота сульфата магния, кГц;} \\
&
\multicolumn{2}{>{\raggedright\arraybackslash}p{(\columnwidth - 4\tabcolsep) * \real{0.9339} + 2\tabcolsep}@{}}{%
\(T\) -- температура воды, °C;} \\
&
\multicolumn{2}{>{\raggedright\arraybackslash}p{(\columnwidth - 4\tabcolsep) * \real{0.9339} + 2\tabcolsep}@{}}{%
\(S\) -- солёность, ‰;} \\
&
\multicolumn{2}{>{\raggedright\arraybackslash}p{(\columnwidth - 4\tabcolsep) * \real{0.9339} + 2\tabcolsep}@{}}{%
\(D\) -- глубина, м;} \\
&
\multicolumn{2}{>{\raggedright\arraybackslash}p{(\columnwidth - 4\tabcolsep) * \real{0.9339} + 2\tabcolsep}@{}}{%
\(pH\) -- водородный показатель воды;} \\
&
\multicolumn{2}{>{\raggedright\arraybackslash}p{(\columnwidth - 4\tabcolsep) * \real{0.9339} + 2\tabcolsep}@{}}{%
\(c\) -- скорость звука в воде, м/с.} \\
\end{longtable}

Выбор формулы осуществляется пользователем.

\section{4.2 Требования к
надёжности}\label{ux442ux440ux435ux431ux43eux432ux430ux43dux438ux44f-ux43a-ux43dux430ux434ux451ux436ux43dux43eux441ux442ux438}

\section{4.2.1 Требования к обеспечению надёжного функционирования
программы}\label{ux442ux440ux435ux431ux43eux432ux430ux43dux438ux44f-ux43a-ux43eux431ux435ux441ux43fux435ux447ux435ux43dux438ux44e-ux43dux430ux434ux451ux436ux43dux43eux433ux43e-ux444ux443ux43dux43aux446ux438ux43eux43dux438ux440ux43eux432ux430ux43dux438ux44f-ux43fux440ux43eux433ux440ux430ux43cux43cux44b}

Некорректные критические параметры конфигурации модуля должны выявляться
с выдачей диагностических сообщений такими средствами ns-3 как механизмы
assert, abort и fatal diagnostics.

\section{4.2.2 Время восстановления после
отказа}\label{ux432ux440ux435ux43cux44f-ux432ux43eux441ux441ux442ux430ux43dux43eux432ux43bux435ux43dux438ux44f-ux43fux43eux441ux43bux435-ux43eux442ux43aux430ux437ux430}

Требования к времени восстановления после отказа не предъявляются, так
как модуль является симуляционной библиотекой и не выполняет
автоматическое восстановление состояния после критических ошибок
конфигурации.

\section{4.2.3 Отказы из-за некорректных действий
пользователя}\label{ux43eux442ux43aux430ux437ux44b-ux438ux437-ux437ux430-ux43dux435ux43aux43eux440ux440ux435ux43aux442ux43dux44bux445-ux434ux435ux439ux441ux442ux432ux438ux439-ux43fux43eux43bux44cux437ux43eux432ux430ux442ux435ux43bux44f}

Для предотвращения отказов из-за некорректных действий пользователя в
документации должны быть описаны допустимые значения ключевых
параметров.

В случае если пользователь запускает программу с некорректными
значениями программа должна завершить выполнение, выдать диагностическое
сообщение и корректно завершить работу.

\section{4.3 Требования к условиям
эксплуатации}\label{ux442ux440ux435ux431ux43eux432ux430ux43dux438ux44f-ux43a-ux443ux441ux43bux43eux432ux438ux44fux43c-ux44dux43aux441ux43fux43bux443ux430ux442ux430ux446ux438ux438}

\section{4.3.1 Климатические условия
эксплуатации}\label{ux43aux43bux438ux43cux430ux442ux438ux447ux435ux441ux43aux438ux435-ux443ux441ux43bux43eux432ux438ux44f-ux44dux43aux441ux43fux43bux443ux430ux442ux430ux446ux438ux438}

Требования к климатическим условиям эксплуатации не предъявляются.

\section{4.3.2 Требования к квалификации и численности
персонала}\label{ux442ux440ux435ux431ux43eux432ux430ux43dux438ux44f-ux43a-ux43aux432ux430ux43bux438ux444ux438ux43aux430ux446ux438ux438-ux438-ux447ux438ux441ux43bux435ux43dux43dux43eux441ux442ux438-ux43fux435ux440ux441ux43eux43dux430ux43bux430}

Для эксплуатации программного модуля достаточно 1 пользователя,
обладающего квалификацией инженера-программиста и владеющего следующими
навыками:

\begin{itemize}
\item
  Навык сборки С++ проектов;
\item
  Базовые навыки работы с библиотекой Ns-3.
\end{itemize}

\section{4.4 Требования к составу и параметрам технических
средств}\label{ux442ux440ux435ux431ux43eux432ux430ux43dux438ux44f-ux43a-ux441ux43eux441ux442ux430ux432ux443-ux438-ux43fux430ux440ux430ux43cux435ux442ux440ux430ux43c-ux442ux435ux445ux43dux438ux447ux435ux441ux43aux438ux445-ux441ux440ux435ux434ux441ux442ux432}

Для работы с программным модулем и библиотекой Ns-3 требуется ПК со
следующими характеристиками:

\begin{itemize}
\item
  Операционная система - Linux;
\item
  Процессор - x86-64, не менее 4 логических ядер;
\item
  ОЗУ - не менее 8 ГБ;
\item
  Свободное место на диске - не менее 5 ГБ (исходный код программы,
  сборка, результаты).
\end{itemize}

\section{4.5 Требования к информационной и программной
совместимости}\label{ux442ux440ux435ux431ux43eux432ux430ux43dux438ux44f-ux43a-ux438ux43dux444ux43eux440ux43cux430ux446ux438ux43eux43dux43dux43eux439-ux438-ux43fux440ux43eux433ux440ux430ux43cux43cux43dux43eux439-ux441ux43eux432ux43cux435ux441ux442ux438ux43cux43eux441ux442ux438}

\section{4.5.1 Требования к информационным структурам и методам
решения}\label{ux442ux440ux435ux431ux43eux432ux430ux43dux438ux44f-ux43a-ux438ux43dux444ux43eux440ux43cux430ux446ux438ux43eux43dux43dux44bux43c-ux441ux442ux440ux443ux43aux442ux443ux440ux430ux43c-ux438-ux43cux435ux442ux43eux434ux430ux43c-ux440ux435ux448ux435ux43dux438ux44f}

Модуль должен поддерживать:

\begin{itemize}
\item
  Механизмы задания параметров конфигурации, принятые в библиотеке Ns-3;
\item
  Исходный код модуля должен быть оформлен в стиле, принятом в
  Ns-3-модулях (структура каталогов, нейминг, документирование).
\end{itemize}

\section{4.5.2 Требования к исходным кодам и языкам
программирования}\label{ux442ux440ux435ux431ux43eux432ux430ux43dux438ux44f-ux43a-ux438ux441ux445ux43eux434ux43dux44bux43c-ux43aux43eux434ux430ux43c-ux438-ux44fux437ux44bux43aux430ux43c-ux43fux440ux43eux433ux440ux430ux43cux43cux438ux440ux43eux432ux430ux43dux438ux44f}

Исходный код должен представлять собой форк репозитория Ns-3 с
разработанным программным модулем.

Компилятор C++ (gcc/clang) с поддержкой стандарта C++20.

\section{4.5.3 Требования к программным средствам, используемым
программой}\label{ux442ux440ux435ux431ux43eux432ux430ux43dux438ux44f-ux43a-ux43fux440ux43eux433ux440ux430ux43cux43cux43dux44bux43c-ux441ux440ux435ux434ux441ux442ux432ux430ux43c-ux438ux441ux43fux43eux43bux44cux437ux443ux435ux43cux44bux43c-ux43fux440ux43eux433ux440ux430ux43cux43cux43eux439}

Для работы модуля используется ns-3 с модулями core, network, mobility,
spectrum, propagation. Для трассировки используется pcap/trace-helper.
Для сборки используется CMake/ns3 build system, компилятор gcc или clang
с поддержкой C++20.

\section{4.5.4 Требования к защите
информации}\label{ux442ux440ux435ux431ux43eux432ux430ux43dux438ux44f-ux43a-ux437ux430ux449ux438ux442ux435-ux438ux43dux444ux43eux440ux43cux430ux446ux438ux438}

Требования к защите информации не предъявляются (модуль не реализует
механизмы аутентификации/авторизации и не обрабатывает персональные
данные).

\section{4.6 Требования к маркировке и
упаковке}\label{ux442ux440ux435ux431ux43eux432ux430ux43dux438ux44f-ux43a-ux43cux430ux440ux43aux438ux440ux43eux432ux43aux435-ux438-ux443ux43fux430ux43aux43eux432ux43aux435}

Требования к маркировке и упаковке не предъявляются. Поставка
программного модуля осуществляется в виде форка библиотеки Ns-3 с
разработанным программным модулем.

\section{4.7 Требования к транспортированию и
хранению}\label{ux442ux440ux435ux431ux43eux432ux430ux43dux438ux44f-ux43a-ux442ux440ux430ux43dux441ux43fux43eux440ux442ux438ux440ux43eux432ux430ux43dux438ux44e-ux438-ux445ux440ux430ux43dux435ux43dux438ux44e}

Хранение исходного кода программы должно быть выполнено с помощью
системы контроля версий Git.

\section{4.8 Специальные
требования}\label{ux441ux43fux435ux446ux438ux430ux43bux44cux43dux44bux435-ux442ux440ux435ux431ux43eux432ux430ux43dux438ux44f}

При разработке программного модуля необходимо использовать стандартные
интерфейсы, используемые в Ns-3 для повышения совместимости с такими
модулями, как network, propagation, mobility, spectrum.

\section{5 Требования к программной
документации}\label{ux442ux440ux435ux431ux43eux432ux430ux43dux438ux44f-ux43a-ux43fux440ux43eux433ux440ux430ux43cux43cux43dux43eux439-ux434ux43eux43aux443ux43cux435ux43dux442ux430ux446ux438ux438}

Состав разрабатываемой программной документации должен включать в себя:

\begin{itemize}
\item
  Техническое задание, оформленное в соответствии с Учебно-методическим
  пособием «Подготовка, оформление выпускной квалификационной работы и
  преддипломной практики»;
\item
  Пояснительную записку, оформленную в соответствии с
  Учебно-методическим пособием «Подготовка, оформление выпускной
  квалификационной работы и преддипломной практики»;
\item
  Руководство системного программиста, оформленное в соответствии с
  учебно-методическим пособием «Подготовка, оформление выпускной
  квалификационной работы и преддипломной практики».
\end{itemize}

\section{6 Технико-экономические
показатели}\label{ux442ux435ux445ux43dux438ux43aux43e-ux44dux43aux43eux43dux43eux43cux438ux447ux435ux441ux43aux438ux435-ux43fux43eux43aux430ux437ux430ux442ux435ux43bux438}

\section{6.1 Экономические преимущества
разработки}\label{ux44dux43aux43eux43dux43eux43cux438ux447ux435ux441ux43aux438ux435-ux43fux440ux435ux438ux43cux443ux449ux435ux441ux442ux432ux430-ux440ux430ux437ux440ux430ux431ux43eux442ux43aux438}

Разработка программного модуля модели программно-аппаратного
гидроакустического приёмопередатчика для проведения виртуальных
испытаний и подготовительных исследований для последующих полунатурных
испытаний обеспечивает сокращение временных и финансовых затрат на
создание и отработку систем гидроакустической связи.

Применение компьютерного моделирования позволяет перенести значительную
часть проверок и исследований с этапа изготовления опытного образца на
этап проектирования. Это снижает объём натурных испытаний, требующих
привлечения специализированного оборудования, испытательных акваторий и
обслуживающего персонала, что приводит к уменьшению прямых
эксплуатационных расходов.

Использование модели на ранних стадиях разработки обеспечивает выявление
архитектурных и алгоритмических ошибок до изготовления аппаратной части.
В результате сокращаются затраты, связанные с внесением изменений в
опытные образцы, переработкой конструкторской документации и повторным
проведением испытаний.

Реализация возможности многократного воспроизведения сценариев
эксплуатации в виртуальной среде позволяет ускорить процесс отработки
алгоритмов функционирования мультиагентных систем и оптимизации
параметров связи. Это способствует сокращению общего срока разработки и
повышению качества проектных решений.

Таким образом, внедрение разработанного программного модуля обеспечивает
снижение совокупной стоимости разработки гидроакустических систем связи,
уменьшение сроков создания опытных образцов и повышение эффективности
проведения испытаний.

\section{7 Стадии и этапы
разработки}\label{ux441ux442ux430ux434ux438ux438-ux438-ux44dux442ux430ux43fux44b-ux440ux430ux437ux440ux430ux431ux43eux442ux43aux438}

\section{7.1 Стадии
разработки}\label{ux441ux442ux430ux434ux438ux438-ux440ux430ux437ux440ux430ux431ux43eux442ux43aux438}

Разработка проходит в шесть этапов в соответствии с таблицей Б.1.

Таблица Б.1 -- Стадии разработки

\begin{longtable}[]{@{}
  >{\raggedright\arraybackslash}p{(\columnwidth - 4\tabcolsep) * \real{0.3194}}
  >{\raggedright\arraybackslash}p{(\columnwidth - 4\tabcolsep) * \real{0.3467}}
  >{\raggedright\arraybackslash}p{(\columnwidth - 4\tabcolsep) * \real{0.3339}}@{}}
\toprule\noalign{}
\endhead
\bottomrule\noalign{}
\endlastfoot
Название этапа & Сроки & Артефакт(-ы) \\
Анализ требований

и предметной области & 05.09.2025 -- 30.09.2025 & Неформализованный
набор требований \\
Формализация

требований & 01.10.2025 -- 31.10.2025 & Техническое задание \\
Проектирование

модуля & 01.11.2025 -- 30.11.2025 & Диаграммы классов,

использования и т.д. \\
Разработка программного модуля & 01.12.2025 -- 28.02.2026 & Рабочий
прототип гидроакустического модуля \\
Тестирование модуля & 01.03.2026 -- 30.04.2026 & Журнал тестирования \\
Документирование & 01.05.2026 -- 31.05.2026 & Пояснительная записка \\
\end{longtable}

\section{7.2 Содержание работ по
этапам}\label{ux441ux43eux434ux435ux440ux436ux430ux43dux438ux435-ux440ux430ux431ux43eux442-ux43fux43e-ux44dux442ux430ux43fux430ux43c}

На стадии анализа требований и предметной области должны быть выполнены
следующие виды работ:

\begin{itemize}
\item
  Сбор требований;
\item
  Исследование предметной области;
\item
  Разработка технического задания.
\end{itemize}

На стадии проектирования модуля должны быть выполнены:

\begin{itemize}
\item
  Разработка диаграммы вариантов использования;
\item
  Разработка диаграммы классов проектного уровня.
\end{itemize}

На стадии разработки модуля должно быть выполнено кодирование модуля.

На стадии тестирования модуля должен быть протестирован функционал
модуля.

На стадии документирования должна быть выполнена разработка
пояснительной записки.

\section{8 Порядок контроля и
приёмки}\label{ux43fux43eux440ux44fux434ux43eux43a-ux43aux43eux43dux442ux440ux43eux43bux44f-ux438-ux43fux440ux438ux451ux43cux43aux438}

\section{8.1 Виды
испытаний}\label{ux432ux438ux434ux44b-ux438ux441ux43fux44bux442ux430ux43dux438ux439}

Программа сдаётся на проверку заказчику 05.06.2026. При обнаружении в
программе ошибок или недостатков исполнитель обязуется устранить их в
недельный срок и предоставить программу на повторную проверку.

Программа сдаётся на проверку независимому тестировщику не позднее
09.06.2026. Результаты тестирования предоставляются на защите дипломного
проекта членам ГАК.

\section{\texorpdfstring{\hfill\break
}{ }}\label{section}

\section{Приложение
Б.1}\label{ux43fux440ux438ux43bux43eux436ux435ux43dux438ux435-ux431.1}

\section{}\label{section-1}

\section{Диаграмма вариантов
использования}\label{ux434ux438ux430ux433ux440ux430ux43cux43cux430-ux432ux430ux440ux438ux430ux43dux442ux43eux432-ux438ux441ux43fux43eux43bux44cux437ux43eux432ux430ux43dux438ux44f}

На рисунке Б.1.1 представлена диаграмма вариантов использования.

Рисунок Б.1.1 -- Диаграмма вариантов использования «Use Case».

\section{Приложение
Б.2}\label{ux43fux440ux438ux43bux43eux436ux435ux43dux438ux435-ux431.2}

\section{Сценарии вариантов
использования}\label{ux441ux446ux435ux43dux430ux440ux438ux438-ux432ux430ux440ux438ux430ux43dux442ux43eux432-ux438ux441ux43fux43eux43bux44cux437ux43eux432ux430ux43dux438ux44f}

1. Запуск симуляции с двумя узлами

1.1 Параметры по умолчанию

\begin{itemize}
\item
  Пользователь запускает сборку тестового примера «передача сообщения
  между двумя узлами с гидроакустическими приёмопередатчиками».
\item
  Пользователь запускает программу без параметров командной строки.
\item
  Если межузловое расстояние не задано, используется значение по
  умолчанию (1 км).
\item
  Программа формирует отчёт симуляции, содержащий:

  \begin{itemize}
  \item
    параметры сценария, включая межузловое расстояние по умолчанию;
  \item
    время отправки сообщения;
  \item
    результат приёма сообщения (успех);
  \item
    время приёма сообщения;
  \item
    адрес устройства-отправителя;
  \item
    адрес устройства-получателя;
  \item
    полезную нагрузку сообщения;
  \item
    среднее значение ОСШ.
  \end{itemize}
\item
  Сообщение принято успешно.
\end{itemize}

1.2 Увеличенное межузловое расстояние

\begin{itemize}
\item
  Пользователь запускает сборку тестового примера «передача сообщения
  между двумя узлами с гидроакустическими приёмопередатчиками».
\item
  Пользователь запускает программу с параметром командной строки,
  задающим межузловое расстояние в 10 км.
\item
  Программа формирует отчёт симуляции, содержащий:

  \begin{itemize}
  \item
    параметры сценария, включая пользовательское значение межузлового
    расстояния;
  \item
    результат приёма сообщения (неудача).
  \end{itemize}
\item
  При достаточно большом межузловом расстоянии сообщение не принимается
  успешно.
\end{itemize}

\hl{2. Запуск симуляции с тремя узлами}

\hl{2.1 Одновременная передача двумя узлами-отправителями при одинаковом
расстоянии до узла-получателя}

\begin{itemize}
\item
  Пользователь запускает сборку тестового примера «одновременная
  передача сообщений двумя узлами-отправителями одному узлу-получателю».
\item
  Пользователь запускает программу без параметров командной строки.
\item
  Если межузловые расстояния не заданы, используются значения по
  умолчанию, при которых оба узла-отправителя находятся на одинаковом
  расстоянии от узла-получателя.
\item
  Программа формирует отчёт симуляции, содержащий:

  \begin{itemize}
  \item
    параметры сценария, включая межузловые расстояния по умолчанию;
  \item
    время отправки первого сообщения;
  \item
    время отправки второго сообщения;
  \item
    результат приёма первого сообщения (неудача);
  \item
    результат приёма второго сообщения (неудача);
  \item
    среднее значение ОСШ для сообщения, если оно доступно;
  \item
    среднее значение ОСШ для сообщения, если оно доступно;
  \end{itemize}
\item
  При одновременной передаче сообщений от двух узлов-отправителей,
  находящихся на одинаковом расстоянии от узла-получателя, оба сообщения
  не принимаются успешно из-за взаимной интерференции.
\end{itemize}

\hl{}

\hl{2.2 Одновременная передача двумя узлами-отправителями при разном
расстоянии до узла-получателя}

\begin{itemize}
\item
  Пользователь запускает сборку тестового примера «одновременная
  передача сообщений двумя узлами-отправителями одному узлу-получателю».
\item
  Пользователь запускает программу с параметрами командной строки,
  задающими межузловые расстояния от каждого узла-отправителя до
  узла-получателя.
\item
  Программа формирует отчёт симуляции, содержащий:

  \begin{itemize}
  \item
    параметры сценария, включая пользовательские значения межузловых
    расстояний;
  \item
    время отправки первого сообщения;
  \item
    время отправки второго сообщения;
  \item
    результат приёма первого сообщения (успех);
  \item
    время приёма первого сообщения;
  \end{itemize}

  \begin{itemize}
  \item
    среднее значение ОСШ для первого сообщения, если оно доступно;
  \end{itemize}

  \begin{itemize}
  \item
    результат приёма второго сообщения (неудача);
  \item
    среднее значение ОСШ для второго сообщения, если оно доступно.
  \end{itemize}
\item
  При одновременной передаче сообщений от двух узлов-отправителей,
  находящихся на разном расстоянии от узла-получателя, сообщение от
  ближнего узла принимается успешно, а сообщение от дальнего узла не
  принимается успешно.\hl{}
\end{itemize}

\section{Приложение
Б.3}\label{ux43fux440ux438ux43bux43eux436ux435ux43dux438ux435-ux431.3}

\section{}\label{section-2}

\section{Макеты экранных
форм}\label{ux43cux430ux43aux435ux442ux44b-ux44dux43aux440ux430ux43dux43dux44bux445-ux444ux43eux440ux43c}

Макеты экранных форм работы модуля представлены на рисунках Б.3.1 -
Б.3.3.

Рисунок Б.3.1 -- Отображение результатов трассировки в программе
Wireshark

Рисунок Б.3.2 -- Отображение результатов работы сценария с двумя
узлами\hl{}

\hl{}

\hl{}

Рисунок Б.3.3 -- Отображение результатов работы сценария с тремя
узлами\hl{}

\section{Приложение
Б.4}\label{ux43fux440ux438ux43bux43eux436ux435ux43dux438ux435-ux431.4}

\section{}\label{section-3}

\section{Структура и формат
данных}\label{ux441ux442ux440ux443ux43aux442ux443ux440ux430-ux438-ux444ux43eux440ux43cux430ux442-ux434ux430ux43dux43dux44bux445}

\hl{}

\section{Б.4.1 Сценарий №1 --- взаимодействие двух
узлов}\label{ux431.4.1-ux441ux446ux435ux43dux430ux440ux438ux439-1-ux432ux437ux430ux438ux43cux43eux434ux435ux439ux441ux442ux432ux438ux435-ux434ux432ux443ux445-ux443ux437ux43bux43eux432}

Б.4.1.1 Входные данные\hl{}

Входные данные представляют собой параметры командной строки \hl{и}
представлены в таблице Б.4.1\hl{}

Таблица Б.4.1 --- Входные данные сценария взаимодействия двух узлов\hl{}

\begin{longtable}[]{@{}
  >{\raggedright\arraybackslash}p{(\columnwidth - 8\tabcolsep) * \real{0.1969}}
  >{\raggedright\arraybackslash}p{(\columnwidth - 8\tabcolsep) * \real{0.2423}}
  >{\raggedright\arraybackslash}p{(\columnwidth - 8\tabcolsep) * \real{0.2120}}
  >{\raggedright\arraybackslash}p{(\columnwidth - 8\tabcolsep) * \real{0.1666}}
  >{\raggedright\arraybackslash}p{(\columnwidth - 8\tabcolsep) * \real{0.1822}}@{}}
\toprule\noalign{}
\begin{minipage}[b]{\linewidth}\raggedright
Параметр командной строки
\end{minipage} & \begin{minipage}[b]{\linewidth}\raggedright
Назначение
\end{minipage} & \begin{minipage}[b]{\linewidth}\raggedright
Тип данных
\end{minipage} & \begin{minipage}[b]{\linewidth}\raggedright
Единица измерения
\end{minipage} & \begin{minipage}[b]{\linewidth}\raggedright
Значение по умолчанию
\end{minipage} \\
\midrule\noalign{}
\endhead
\bottomrule\noalign{}
\endlastfoot
-\/-sound-speed & Скорость распространения звука в среде & вещественное
число & м/с & 1500 \\
-\/-geometric-coefficient & Коэффициент геометрического расхождения &
вещественное число & нет & 2.0 \\
-\/-absorption-model & Имя модели поглощения сигнала & строка & нет &
ns3::FrancoisGarrisonPropagationLossModel \\
-\/-snr-threshold & Минимальное среднее значение ОСШ, необходимое для
успешного приёма & вещественное число & нет & 10 \\
-\/-distance & Межузловое расстояние между узлом-отправителем и
узлом-получателем & вещественное число & м & 1000 \\
\end{longtable}

Б.4.1.2 Выходные данные\hl{}

\hl{Выходные данные формируются в текстовом виде и выводятся в
стандартный поток вывода (stdout) и имеют следующий вид:}

\emph{\hl{Simulation with 2 nodes.}}

\emph{\hl{Parameters:}}

\emph{\hl{Inter-node distance: \textless значение\textgreater{} m}}

\emph{\hl{Sound speed: \textless значение\textgreater{} m/s}}

\emph{\hl{Geometric spreading loss coefficient:
\textless значение\textgreater{}}}

\emph{\hl{Absorption loss model: \textless значение\textgreater{}}}

\emph{\hl{SNR threshold: \textless значение\textgreater{}}}

\emph{\hl{Carrier frequency: \textless значение\textgreater{} Hz}}

\emph{\hl{Bandwidth: \textless значение\textgreater{} Hz}}

\emph{\hl{Power: \textless значение\textgreater{} uPa\^{}2}}

\emph{\hl{PSD: \textless значение\textgreater{} uPa\^{}2/Hz}}

\emph{\hl{Send time: \textless значение\textgreater{} ms}}

\emph{\hl{Receive result: success\textbar failure}}

\hl{При успешном приёме дополнительно выводятся:}

\emph{\hl{Receive time: \textless значение\textgreater{} ms}}

\emph{\hl{Sender address: \textless значение\textgreater{}}}

\emph{\hl{Receiver address: \textless значение\textgreater{}}}

\emph{\hl{Payload: \textless значение\textgreater{}}}

\emph{\hl{Mean SNR: \textless значение\textgreater{}}}

\hl{При неуспешном приёме дополнительно выводится:}

\hl{\emph{Message was not received successfully.}}

\section{\texorpdfstring{\hl{}}{}}\label{section-4}

\section{Б.4.2 Сценарий №2 --- взаимодействие трёх
узлов}\label{ux431.4.2-ux441ux446ux435ux43dux430ux440ux438ux439-2-ux432ux437ux430ux438ux43cux43eux434ux435ux439ux441ux442ux432ux438ux435-ux442ux440ux451ux445-ux443ux437ux43bux43eux432}

Б.4.2.1 Входные данные\hl{}

Входные данные представляют собой параметры командной строки \hl{и}
представлены в таблице Б.4.2\hl{}

Таблица Б.4.2 --- Входные данные сценария взаимодействия трёх узлов\hl{}

\begin{longtable}[]{@{}
  >{\raggedright\arraybackslash}p{(\columnwidth - 8\tabcolsep) * \real{0.1969}}
  >{\raggedright\arraybackslash}p{(\columnwidth - 8\tabcolsep) * \real{0.2423}}
  >{\raggedright\arraybackslash}p{(\columnwidth - 8\tabcolsep) * \real{0.2120}}
  >{\raggedright\arraybackslash}p{(\columnwidth - 8\tabcolsep) * \real{0.1666}}
  >{\raggedright\arraybackslash}p{(\columnwidth - 8\tabcolsep) * \real{0.1822}}@{}}
\toprule\noalign{}
\begin{minipage}[b]{\linewidth}\raggedright
Параметр командной строки
\end{minipage} & \begin{minipage}[b]{\linewidth}\raggedright
Назначение
\end{minipage} & \begin{minipage}[b]{\linewidth}\raggedright
Тип данных
\end{minipage} & \begin{minipage}[b]{\linewidth}\raggedright
Единица измерения
\end{minipage} & \begin{minipage}[b]{\linewidth}\raggedright
Значение по умолчанию
\end{minipage} \\
\midrule\noalign{}
\endhead
\bottomrule\noalign{}
\endlastfoot
-\/-sound-speed & Скорость распространения звука в среде & вещественное
число & м/с & 1500 \\
-\/-geometric-coefficient & Коэффициент геометрического расхождения &
вещественное число & нет & 2.0 \\
-\/-absorption-model & Имя модели поглощения сигнала & строка & нет &
ns3::FrancoisGarrisonPropagationLossModel \\
-\/-snr-threshold & Минимальное среднее значение ОСШ, необходимое для
успешного приёма & вещественное число & нет & 10 \\
-\/-sender1-distance & Расстояние от первого узла-отправителя до
узла-получателя & вещественное число & м & 1000 \\
-\/-sender2-distance & Расстояние от второго узла-отправителя до
узла-получателя & вещественное число & м & 1000 \\
\end{longtable}

Б.4.2.2 Выходные данные\hl{}

\hl{Выходные данные формируются в текстовом виде и выводятся в
стандартный поток вывода (stdout) и имеют следующий вид:}

\emph{\hl{Simulation with 3 nodes.}}

\emph{\hl{Parameters:}}

\emph{\hl{Sender 1 distance: \textless значение\textgreater{} m}}

\emph{\hl{Sender 2 distance: \textless значение\textgreater{} m}}

\emph{\hl{Sound speed: \textless значение\textgreater{} m/s}}

\emph{\hl{Geometric spreading loss coefficient:
\textless значение\textgreater{}}}

\emph{\hl{Absorption loss model: \textless значение\textgreater{}}}

\emph{\hl{SNR threshold: \textless значение\textgreater{}}}

\emph{\hl{Carrier frequency: \textless значение\textgreater{} Hz}}

\emph{\hl{Bandwidth: \textless значение\textgreater{} Hz}}

\emph{\hl{Power: \textless значение\textgreater{} uPa\^{}2}}

\emph{\hl{PSD: \textless значение\textgreater{} uPa\^{}2/Hz}}

\emph{\hl{Message 1:}}

\emph{\hl{Send time: \textless значение\textgreater{} ms}}

\emph{\hl{Receive result: success\textbar failure}}

\emph{\hl{Mean SNR: \textless значение\textgreater\textbar not
available}}

\emph{\hl{Message 2:}}

\emph{\hl{Send time: \textless значение\textgreater{} ms}}

\emph{\hl{Receive result: success\textbar failure}}

\emph{\hl{Mean SNR: \textless значение\textgreater\textbar not
available}}

\hl{При успешном приёме для каждого сообщения дополнительно выводятся:}

\emph{\hl{Receive time: \textless значение\textgreater{} ms}}

\emph{\hl{Sender address: \textless значение\textgreater{}}}

\emph{\hl{Receiver address: \textless значение\textgreater{}}}

\emph{\hl{Payload: \textless значение\textgreater{}}}\\
\hl{}

\section{Приложение
Б.5}\label{ux43fux440ux438ux43bux43eux436ux435ux43dux438ux435-ux431.5}

\section{}\label{section-5}

\section{Состояния физического уровня модели
приёмопередатчика}\label{ux441ux43eux441ux442ux43eux44fux43dux438ux44f-ux444ux438ux437ux438ux447ux435ux441ux43aux43eux433ux43e-ux443ux440ux43eux432ux43dux44f-ux43cux43eux434ux435ux43bux438-ux43fux440ux438ux451ux43cux43eux43fux435ux440ux435ux434ux430ux442ux447ux438ux43aux430}

На рисунке Б.5.1 представлена диаграмма машины состояний.

\section{}\label{section-6}

Рисунок Б.5.1 --- Диаграмма состояний физического уровня модели
приёмопередатчика
